\chapter{结果分析与讨论}\label{chap:analysis}

\section{不同算法的表现差异}

从实验结果可以看出，BFS 与 Dijkstra 都能够稳定找到最短路径，但在复杂场景中访问节点数很多，说明这两类算法虽然可靠，但搜索效率偏低。A* 在引入启发式函数后，访问节点明显减少，说明启发式信息确实能够有效缩小搜索范围。

改进 A* 的路径长度与普通 A* 一致，但综合代价更低，这说明加入贴障惩罚和转向惩罚后，路径质量得到了提升。虽然它的访问节点数不一定最少，但在“更安全、更平滑”这一目标上表现更好。

\section{Dijkstra 与其他算法比较}

在本文的实验中，Dijkstra 是一个很有代表性的对照算法。由于当前栅格地图中每一步的基础移动代价相同，所以 Dijkstra 与 BFS 在“是否能找到最短路径”这一点上表现接近，二者都能够稳定得到最短长度路径。但从搜索过程来看，Dijkstra 仍然需要像 BFS 一样扩展大量节点，因此在复杂场景中的效率不高。

与 A* 相比，Dijkstra 最大的区别在于它不使用终点方向信息。A* 会借助启发式函数优先扩展更接近终点的节点，所以在 30x30 场景中访问节点数明显少于 Dijkstra。这说明在静态已知地图中，启发式信息确实能够有效降低搜索范围。

与改进 A* 相比，Dijkstra 的优点是思路简单、最短路性质明确，不需要额外调参；缺点是无法表达“贴障风险”和“转弯代价”这类路径质量信息。改进 A* 虽然不一定在节点数上绝对占优，但生成的路径通常更平滑，也更符合机器人实际运动需求。

与 Q-Learning 相比，Dijkstra 不需要训练过程，结果更稳定；但它缺少学习能力，不能像强化学习方法那样通过与环境交互逐步形成策略。与人工势场法相比，Dijkstra 不会因为局部极小值问题而失败，因此在求解稳定性上更有优势。

总体来看，Dijkstra 适合作为一个“可靠但效率一般”的基准算法。在本次课程设计中，它的意义主要是提供一个无启发式最短路参考，从而更清楚地看出 A*、改进 A* 和 Q-Learning 等方法各自带来的改进效果。

\section{强化学习与启发式方法的讨论}

Q-Learning 在训练足够轮次后能够成功找到路径，说明强化学习方法在此类问题上具有可行性。相比图搜索算法，它的优势在于可以通过交互学习策略；不足之处在于训练时间较长，而且结果对奖励函数设计比较敏感。

人工势场法在简单环境中表现较直观，但在本文设计的复杂 30x30 场景中出现了失败，原因主要是 U 型障碍和窄通道导致局部极小值问题。这也说明人工势场法更适合作为局部避障方法，而不太适合单独承担复杂场景下的全局规划任务。

\section{当前方案的局限性}

虽然本文已经完成了多种算法的实现与比较，但仍存在一些局限：

\begin{itemize}
  \item 当前地图是静态已知环境，没有加入动态障碍；
  \item 机器人运动方式只考虑上下左右四个方向，没有考虑速度和转向半径等真实约束；
  \item Q-Learning 仍然是表格型实现，状态规模扩大后训练成本会快速增加；
  \item 动态展示目前以 GIF 形式导出，交互性还不够强。
\end{itemize}

\section{可改进方向}

后续可以从以下几个方向继续改进：

\begin{enumerate}
  \item 引入 D* Lite、LPA* 等增量式搜索算法，处理动态环境变化；
  \item 将 Q-Learning 扩展为深度强化学习方法，提升大规模场景下的泛化能力；
  \item 在代价函数中进一步加入转弯半径、安全距离等约束，使路径更贴近真实机器人需求；
  \item 设计更丰富的可视化界面，实现交互式播放和参数调节。
\end{enumerate}
